\chapter{Obtaining Textual Data}\label{obtaining-textual-data}

\textbf{Learning Objectives} After reading this chapter, you should be
able to:

\emph{Overview starts}

. Understanding where to look for texts to analyze\\
. Downloading datasets from the internet and import them in R\\
. Collecting unstructured data and put them into a structured form\\
. Scraping data with APIs\\
. Scraping data with a web browser's plugins\\
. Programming a script that downloads data for you\\
. Understanding other ways of obtaining text

\emph{Overview ends}

\section{Handling electronic text}\label{handling-electronic-text}

\subsection{Introduction}\label{introduction-1}

Text mining and analytics provide marketing analysts with various tools for
analyzing large volumes of textual data (Guyt, Datta, and Boegershausen
2024). These applications and the important role they play were 
outlined in the introduction. In this section, we will focus on
the technical aspects of text acquisition, which typically takes a 
considerable length of time, and it is important to acknowledge that data
collection and preparation represent a significant portion of a
project's duration, often up to 80\%. Efficiently managing these
processes can significantly reduce the frustration and inefficiency inherent in the acquisition of textual data (Frye
and Schmitt 2020).

A large proportion of data available on the internet is unstructured,
(cf., Gandomi and Haider 2015; Wedel and Kannan 2016; and Boegershausen
et al. 2022). This means it is not organized in a tabular format or
within a database, and lacks a predefined structure. The primary types of
data requiring transformation are text and images. Since computers
operate more effectively with structured data, converting them from
unstructured to structured formats is critical in the analytical
process. Therefore, before commencing data collection, it is essential
to consider several fundamental questions. The foremost concerns the format of the data to be collected. This leads to the
realization that data, in our case text, can be stored in various file
types, such as comma separated value (CSV) files, which are appropriate for
this task. Additionally, it is good to understand the relationships
within the data, the number of sources, and their acquisition methods.
We speak using words, but it is through written text that we leave
our imprint on the future.

Humans continuously generate text in various forms, ranging from
physical mediums such as letters, articles, diaries, books, and other
written formats, to virtual ones stored in file types such as
.\textbf{txt}, .\textbf{doc}, .\textbf{docx}, .\textbf{pages},
.\textbf{pdf}, and .\textbf{odt} files, among others. Additionally, there are file
formats with more or less familiar extensions. In text analysis or
analytics, having documents in an electronic format is imperative, as it
allows the text to be structured and manipulated easily. Consequently, acquiring
virtual text is preferable, but there are methods to convert printed
text into electronic documents. Electronic text can be derived from
printed material through a two-step process. The initial step involves
digitization, achieved by scanning or photographing the documents. The
subsequent step is to extract text from these images, using
optical character recognition (OCR) systems. This process is often
lengthy and prone to errors (Srihari 1985), requiring specific
software and human intervention for quality assurance. Proofreading is
critical, requiring meticulous attention to detail to ensure the transformed text
accurately reflects the original. This procedure applies to
standard printed text, but the challenge intensifies with handwritten
documents.

Mass digitization projects have been, and continue to be, significant
endeavors. Official initiatives like Europeana, backed by the European
Union, have created digital platforms housing millions of items. Another notable project is that of
Google Books, launched with the ambitious goal of digitizing
the world's books and making them accessible online (Thylstrup 2019). A
significant aspect of this project was the use of OCR technology, which converted the scanned images of
books into machine-readable text. However, this project faced legal
challenges from authors and publishers until 2016, when Google won a series of important legal battles and it was allowed to proceed with
its scanning efforts (Rimmer 2017). Today, a number of projects and
partnerships give research groups access to this extensive database
to carry out in-depth text analysis.

For those not requiring access to rare or non-digital textual sources,
using existing electronic versions is a practical approach.
However, although there are official repositories, shadow libraries do exist and it is
important to be mindful of the copyright laws in your region, as not all
digitized texts are freely available for research or other purposes.

Some ways to obtain electronic text are the following: (i)
self-generated text, (ii) manual copy and paste, (iii) direct text
downloads, (iv) web scraping, and (v) utilizing ready-to-use datasets. These
are among the most common approaches, see Figure~\ref{fig-textdata}.
Manual copy and paste is a straightforward yet time-consuming method,
ideal for small amounts of text. Directly downloading text files, such
as in \textbf{pdf} format, \textbf{docx} format, or plain text files, is
efficient for accessing large documents or books already in a digital
format. Scraping text from documents with various file extensions often
requires specialized software or scripts, especially when handling large
volumes of data or complex formats. On the other hand, web scraping ranges from the
simple extraction of specific pieces of text from web pages to more
advanced scraping that navigates web pages and handles dynamic
content. Lastly, ready-to-use datasets from online
repositories are a convenient source, as these are often
well-structured and ideal for analysis without extensive preprocessing.
Next we describe these common ways to obtain textual
data.

\begin{figure}

\centering{

\includegraphics[width=1\textwidth,height=\textheight]{figure/chapter3/obtain_textual_data.png}

}

\caption{\label{fig-textdata}Obtain textual data \newline Authors' own
figure}

\end{figure}%

\subsection{Copy and paste}\label{copy-and-paste}

The simplest method of obtaining text that is not self-generated is
through copy-and-paste. This involves copying text from
various sources and pasting it into a simple text file stored locally on
your computer. It is a three-step process: selecting the text, mainly
with the help of your mouse, copying it, and then pasting it into a
recipient file. You can paste different portions of text consecutively 
onto the same text file, making this procedure a primitive form of
scraping. The procedure is similar in all operating systems, such as
MacOS, Windows, or Linux, with minor differences in the commands and
text editing applications used. Here you have the procedure detailed for
the most common commercial operating systems:

\emph{First, open the source document or webpage from which you wish to
copy text.\\
}Click and drag your mouse over the text to highlight it.\\
\emph{On OSX, press Command + C to copy the highlighted text; on
Windows, use Ctrl + C.\\
}Next, open a text editing application. For MacOS, use TextEdit; for
Windows, use Notepad. Note that TextEdit defaults to rich text mode. To
save your file as a .txt file, convert it to plain text by going to
Format in the menu bar and selecting Make Plain Text (or use the
shortcut Shift + Command + T).\\
\emph{Click inside the document and press Command + V on MacOS, or Ctrl + V
on Windows, to paste the copied text.\\
}To save the file, click File \textgreater{} Save (or use Command + S on
MacOS, or Ctrl + S on Windows). Name your file, for instance,
``copied\_text.txt''. Navigate to the desired folder within your working
directory. If you're looking for a specific folder such as `data' and it
does not exist, create it. Finally, click Save.

This method is straightforward and does not require specialized
software, making it an excellent starting point for anyone new to data
collection.

This is how you import the file into R and visualize the text.

\begin{Shaded}
\begin{Highlighting}[]
\CommentTok{\# Read the text file from the \textquotesingle{}data\textquotesingle{} folder}
\NormalTok{text\_data }\OtherTok{\textless{}{-}} \FunctionTok{readLines}\NormalTok{(}\StringTok{"data/copied\_text.txt"}\NormalTok{)}

\CommentTok{\# Display the content of the text file}
\FunctionTok{print}\NormalTok{(text\_data)}
\end{Highlighting}
\end{Shaded}

\subsection{Downloaded text}\label{downloaded-text}

For a marketing practitioner, sourcing documents which are ready for analysis is
not always straightforward. For example, literature students have the
advantage of accessing many out-of-copyright texts available on the
internet. Alongside well-known resources like Google Books and
Europeana, they (or anyone else) can explore sites like the Gutenberg Project,
Internet Archive, Open Culture, Hathi Trust, ACLS Humanities E-Book,
National Digital Library of India, Gallica, Wikisource, Runiverse, as well as
many others. These platforms offer extensive collections of books and
freely accessible texts, ideal for downloading and importing into R. The
benefit of using simple text is that it allows for imports from these
websites, providing an excellent opportunity to start developing text
analytics skills. The goal is to have the text in a format that lends itself
to structured analysis.

Fortunately, collections of documents are already prepared in structured
formats such as CSV or zipped CSV files. These are often found on platforms
like Kaggle, GitHub, and the repositories of various research institutes. When
dealing with documents in compressed formats like \textbf{.gz} or
\textbf{.zip}, it is advisable to decompress them first before importing
them into R. Conversely, files in formats like \textbf{.csv} or
\textbf{.xlsx} can be imported directly from the internet via a secure
HTTPS connection. Downloading files over an insecure HTTP protocol is
not recommended for security reasons. Direct-internet import can be
cumbersome and prone to errors for files with substantial amounts of data. Hence,
it is preferable to download these files to your local machine first and
then import them into R.

Let us demonstrate this by  importing a dataset directly. This method is suitable for smaller files: Downloading the files first 
is recommended for larger ones. The ``sample.csv'' file, an excerpt from
the larger ``Product Reviews.csv,'' contains the first 250 observations.
This file is an example of a structured document. As shown in
Table~\ref{tbl-csvtable}, the first row, which is the header, lists the
column names or features. Each subsequent row represents an observation,
encompassing all the data related to a single purchase.

\begin{longtable}[]{@{}lllll@{}}
\caption{Table Sample, Excel-like}\label{tbl-csvtable}\tabularnewline
\toprule\noalign{}
Review ID & Item ID & Base Item ID & Review Date & \ldots{} \\
\midrule\noalign{}
\endfirsthead
\toprule\noalign{}
Review ID & Item ID & Base Item ID & Review Date & \ldots{} \\
\midrule\noalign{}
\endhead
\bottomrule\noalign{}
\endlastfoot
R100E6MT94PK6L & B0051VVOB2 & & 1/8/2012 & \ldots{} \\
R100HU42LKLLD0 & B0057O9O6K & & 4/10/2012 & \ldots{} \\
R100PBXP49AU0B & B0051VVOB2 & & 3/26/2012 & \ldots{} \\
\ldots{} & \ldots{} & \ldots{} & \ldots{} & \ldots{} \\
\end{longtable}

The \texttt{read.csv} function in R is used to import the
\texttt{sample.csv} file and automatically transform it into a
dataframe. Remember, in R, a dataframe is a term used to describe a
tabular data structure. The parameter \texttt{header\ =\ TRUE} within
the \texttt{read.csv} function specifies that the first row of the CSV
file should be treated as the header containing the column names rather
than regular data.

\emph{Program code starts}

\begin{Shaded}
\begin{Highlighting}[]
\CommentTok{\# URL of the file}
\NormalTok{url }\OtherTok{\textless{}{-}} \StringTok{"https://danieldanlab.github.io/data/sample.csv"}

\CommentTok{\# Read CSV into dataframe}
\NormalTok{df }\OtherTok{\textless{}{-}} \FunctionTok{read.csv}\NormalTok{(url, }\AttributeTok{header =} \ConstantTok{TRUE}\NormalTok{)}

\CommentTok{\# Check dataset in viewer, uncomment}
\CommentTok{\# View(df)}

\CommentTok{\# Keep only first 5 columns}
\NormalTok{df }\OtherTok{\textless{}{-}}\NormalTok{ df[, }\DecValTok{1}\SpecialCharTok{:}\DecValTok{5}\NormalTok{]}

\CommentTok{\# Display first 6 rows}
\FunctionTok{head}\NormalTok{(df)}
\end{Highlighting}
\end{Shaded}

\begin{verbatim}
  X      Review_ID    Item_ID Base_item_ID Review_date
1 1 R100E6MT94PK6L B0051VVOB2                 1/8/2012
2 2 R100HU42LKLLD0 B0057O9O6K                4/10/2012
3 3 R100PBXP49AU0B B0051VVOB2                3/26/2012
4 4 R100U020EZR95F B0051VVOB2                6/14/2012
5 5 R101ANBUYVZCUM B004U9USEA                5/21/2011
6 6 R101MUGRC35UWS B001UHPRSW                1/29/2010
\end{verbatim}

\begin{Shaded}
\begin{Highlighting}[]
\CommentTok{\# Display last 6 rows}
\FunctionTok{tail}\NormalTok{(df)}
\end{Highlighting}
\end{Shaded}

\begin{verbatim}
      X      Review_ID    Item_ID Base_item_ID Review_date
245 245 R10VM77VD3D76X B004Q1P0S6               12/25/2011
246 246 R10VXS69HI0F1X B0044UAQMY                6/21/2011
247 247 R10W6FDN5GFTK4 B004U78J1G   B004U78J1G  10/24/2011
248 248 R10WG1DINHPFRU B0051VVOB2               12/12/2011
249 249 R10WJVG8HFG133 B0051VVOB2               11/30/2011
250 250  R10WL5UATT2QP B004PGMFG2   B0067XX3AI   6/30/2011
\end{verbatim}

\emph{Program code ends}

The \texttt{df} dataframe can be sliced both horizontally and vertically
to suit your needs. The general syntax for this operation is
\texttt{df{[}rows,\ columns{]}}. In this structure, the brackets
encompass the entire dataframe, and the comma separates the rows from
the columns, in specific order. To select the first row along with
all its columns, you would use d\texttt{f{[}1,\ {]}}, where the (empty)
space after the comma signifies the inclusion of all the columns. Similarly,
to select specific rows and columns, for example, rows 3 to 5 and
columns 2 to 9, you would use \texttt{df{[}3:5,\ 2:9{]}}.

\subsubsection{Common file-type
extensions}\label{common-file-type-extensions}

What if the file you need to analyze does not have a \texttt{.csv}
extension? In such instances, the \texttt{readtext} package (Benoit and
Obeng 2020) is handy. It is designed to read and extract text from a
variety of file formats, including \textbf{.doc}, \textbf{.docx},
\textbf{.html}, \textbf{.json}, \textbf{.pdf}, \textbf{.rtf},
\textbf{.tab}, \textbf{.xls}, \textbf{.xlsx}, \textbf{.xml}, and several
other related extensions. This versatile function not only extracts the
text but also retrieves the document's metadata, and it can handle
multiple file types simultaneously. The data extracted is stored in a
dataframe, our preferred tabular data format. As a side note, when
importing data into R, it is not uncommon to encounter issues with
character encoding. In such cases, check the solutions suggested in the
introduction.

The \texttt{readtext} package relies on other packages for specific file
types: \texttt{pdftools} (Ooms 2020) for \textbf{.pdf} files and
\texttt{antiword} for \textbf{.doc} documents. Installation on Windows
and MacOS is typically straightforward. However, Linux users might face
challenges, particularly with \texttt{pdftools}, which requires a
development library. On Debian/Ubuntu and similar systems, this issue
can be resolved by installing the \texttt{libpoppler} development
library. Other Linux distributions with different package management
systems, like those using \texttt{rpm} package management, have
similar solutions.

\begin{Shaded}
\begin{Highlighting}[]
\CommentTok{\# Installs the libpoppler{-}cpp{-}dev library on Debian/Ubuntu systems}
\FunctionTok{sudo}\NormalTok{ apt install libpoppler{-}cpp{-}dev}
\end{Highlighting}
\end{Shaded}

Text files are often unstructured. This means
that the data are not stored in a tabular form, making them difficult to
work with. To illustrate this, in the next example we have created a
\textbf{.pdf} file that simulates a real-world scenario. This file
consists of two pages. The first page contains text divided into
paragraphs, generated using Lorem Ipsum (2024), a
commonly used placeholder text which imitates Latin. The second page features a table filled with
randomly generated data (Random Data Generator 2024). While this
text appears structured to our eyes, this is not the case for a computer. We
will show how to extract and format this information into a structure
suitable for basic text analytics. The goal is to organize the data into
a simple table, with one column for each sentence and another for an
incremental number to track the paragraphs, similar to the structure
shown in Table~\ref{tbl-csvtable}. It is important to know, a
\textbf{.pdf} file is a container rather than a specific format. If the
document contains text in an image format instead of selectable text,
you should use the \texttt{ocr\ =\ TRUE} option with the
\texttt{pdf\_text()} function to read it.

Let us now download the
file and extract the text from the first page of the \textbf{.pdf}. To
understand its content and structure, it is advisable to open and
inspect the PDF file, as well as any other text file, before working with it.

\begin{Shaded}
\begin{Highlighting}[]
\CommentTok{\# Load necessary libraries}
\FunctionTok{library}\NormalTok{(pdftools)}
\FunctionTok{library}\NormalTok{(stringr)}

\CommentTok{\# Define URL of the PDF file}
\NormalTok{url }\OtherTok{\textless{}{-}} \StringTok{"https://danieldanlab.github.io/data/sample.pdf"}

\CommentTok{\# Download PDF file}
\FunctionTok{download.file}\NormalTok{(url, }\StringTok{"sample.pdf"}\NormalTok{, }\AttributeTok{mode =} \StringTok{"wb"}\NormalTok{)}

\CommentTok{\# Read PDF; each page becomes an element in \textquotesingle{}text\textquotesingle{}}
\NormalTok{text }\OtherTok{\textless{}{-}} \FunctionTok{pdf\_text}\NormalTok{(}\StringTok{"sample.pdf"}\NormalTok{)}

\CommentTok{\# Remove \textquotesingle{}Sample text\textquotesingle{} from the text}
\NormalTok{text }\OtherTok{\textless{}{-}} \FunctionTok{str\_replace}\NormalTok{(text, }\StringTok{"Sample text"}\NormalTok{, }\StringTok{""}\NormalTok{)}

\CommentTok{\# Replace newline characters with space in page 1}
\NormalTok{text[}\DecValTok{1}\NormalTok{] }\OtherTok{\textless{}{-}} \FunctionTok{str\_replace\_all}\NormalTok{(text[}\DecValTok{1}\NormalTok{], }\StringTok{"}\SpecialCharTok{\textbackslash{}n}\StringTok{"}\NormalTok{, }\StringTok{" "}\NormalTok{)}

\CommentTok{\# Split text at periods in page 1}
\NormalTok{text[}\DecValTok{1}\NormalTok{] }\OtherTok{\textless{}{-}} \FunctionTok{str\_split}\NormalTok{(text[}\DecValTok{1}\NormalTok{], }\StringTok{"}\SpecialCharTok{\textbackslash{}\textbackslash{}}\StringTok{."}\NormalTok{)}

\CommentTok{\# Convert to dataframe}
\NormalTok{mytext }\OtherTok{\textless{}{-}} \FunctionTok{as.data.frame}\NormalTok{(text[}\DecValTok{1}\NormalTok{])}

\CommentTok{\# Name the column as \textquotesingle{}Sentence\textquotesingle{}}
\FunctionTok{colnames}\NormalTok{(mytext) }\OtherTok{\textless{}{-}} \FunctionTok{c}\NormalTok{(}\StringTok{"Sentence"}\NormalTok{)}

\CommentTok{\# Display first six rows}
\FunctionTok{head}\NormalTok{(mytext)}
\end{Highlighting}
\end{Shaded}

\begin{verbatim}
                                                       Sentence
1       Lorem ipsum dolor sit amet, consectetur adipiscing elit
2                                      Quisque sed venenatis mi
3  Mauris eget lectus luctus, molestie enim nec, tincidunt urna
4                              Nulla iaculis ut lorem a blandit
5             Etiam non arcu et magna tempor iaculis ac at erat
6          Quisque non malesuada augue, sit amet vulputate enim
\end{verbatim}

Having achieved the primary goal of extracting the data, you can
refine the appearance of the content in the dataset further. This can be done
using regular expressions and string manipulation functions to tailor
the data to your specific needs.

The table on the second page of the \textbf{.pdf} file is already in a
tabular format, which we will aim to keep. To accomplish this, we will
use the \texttt{extract\_tables()} function, which efficiently extracts
the table's content and places it in a dataframe. It is worth noting that in scenarios
where a PDF page contains multiple tables, possibly in various
locations, these can also be extracted. For more advanced options and the
handling of such cases, refer to the \texttt{tabulapdf} package, which
offers additional tools for extracting PDF tables.

\begin{Shaded}
\begin{Highlighting}[]
\CommentTok{\# Load tabulapdf library}
\FunctionTok{library}\NormalTok{(tabulapdf)}

\CommentTok{\# Extract table from second page of PDF}
\NormalTok{atable }\OtherTok{\textless{}{-}} \FunctionTok{extract\_tables}\NormalTok{(url, }\AttributeTok{pages =} \DecValTok{2}\NormalTok{, }\AttributeTok{guess =} \ConstantTok{TRUE}\NormalTok{, }\AttributeTok{output =} \StringTok{"tibble"}\NormalTok{)}

\CommentTok{\# Assign first table to mytable}
\NormalTok{mytable }\OtherTok{\textless{}{-}}\NormalTok{ atable[[}\DecValTok{1}\NormalTok{]]}

\CommentTok{\# Display top six rows of the table}
\FunctionTok{head}\NormalTok{(mytable)}
\end{Highlighting}
\end{Shaded}

\begin{verbatim}
# A tibble: 6 x 5
  `First Name` `Last Name` Gender Education Occupation   
  <chr>        <chr>       <chr>  <chr>     <chr>        
1 Kate         Jones       Female Bachelor  Scientist    
2 Adele        Barnes      Female Doctoral  Social Worker
3 Sofia        Campbell    Female Primary   Mathematician
4 Oscar        Parker      Male   Master    Mathematician
5 Andrew       Armstrong   Male   Bachelor  Mathematician
6 Sawyer       Taylor      Male   Doctoral  Social Worker
\end{verbatim}

All in all, our primary goal is to extract the most crucial element for
our analysis, which is text, regardless of the file format in which it is encoded. The \textbf{.pdf} file type we just discussed is only one example of a
source from which we can obtain information. As you have learned,
importing text data into R from common file formats is a (relatively)
straightforward process. This procedure is similar to that used when
extracting information from other file types, such as Excel formats
(\textbf{.xls} or \textbf{.xlsx}). There are several packages in R
designed for this purpose. For instance, the \texttt{readxl} package is
an excellent choice for reading Excel files. When it comes to reading
files from the internet, the file's extension guides us to choose the
appropriate function. For \textbf{.dat}, \textbf{.txt}, or \textbf{.csv}
files, we can use the \texttt{read.table()} function. R can also import various file types from other software, such as
SPSS. This is made possible with the foreign package and its
\texttt{read.spss()} function. In the next section, we will
tackle more advanced procedures to structure data from the internet
effectively within R.

\subsection{Scraping for data}\label{scraping-for-data}

It usually requires a more elaborate procedure to collect customer opinions and other user-generated content, as they are typically embedded
within specific web pages. How do we collect reviews about a
specific product or service? Many companies, such as Airbnb, Amazon, Booking.com, TripAdvisor, Trustpilot,
and Yelp, host these data on their websites. Similarly, opinion-based text from social-media platforms like
Twitter, Facebook, and Instagram can be valuable. We can access this
data through a process known as web scraping (Khder 2021). As
the term suggests, this involves extracting data from a surface. In a
digital context, it refers to the technique of copying data embedded in
an unstructured way within HTML tags to a structured format which suits the needs of data analysts, such as a
CSV file. What can we scrape
from web pages? Essentially anything digital: text, numbers, images, and
metadata. For instance, in the context of online product feedback, we
can collect customer reviews, product ratings, customer nicknames,
labels, dates, and other related data or metadata. All this can be
transformed into useful information for analysis and knowledge
extraction. Researchers in a number of fields such as economics
(Cavallo 2012; Hillen 2019), psychology (Landers et al. 2016),
agriculture (Yang, Wilson, and Wang 2010), and the hotel industry (Han and Anderson 2021), as well as geographic data acquisition (Brenning
and Henn 2023), obtain data through scraping. Web scraping for consumer
research (Boegershausen et al. 2022) plays a significant role too. There
are several methods and tools which facilitate the downloading and collection of data. Some of them must be built specifically for
the fields in which they are used (Bonifacio et al. 2015), while others
are more generic, such as Outwit (2024), Fminer
(2024), or Parsehub (2024). Some methods need to be tailored to specific projects and
frequently use computer languages such as Java, Python, or R to carry out their
data crawling. These tools change data from an unstructured into
a structured form, making them ready for further analysis. It is important to
note that data, mainly text data, are often the core business of some
companies, and they may try to block or hinder automatic parsers or
robots from scraping and collecting information (Hirschey 2014; Poggi et
al. 2007). Technically, there are two main ways to scrape internet data:
an application programming interface (API), for example {[}sturm2017{]},
or document-object-model (DOM) parsing (Radilova et al. 2022). APIs,
provided by platforms such as Twitter, Facebook, and TripAdvisor, offer
well established procedures for obtaining data in predefined formats such as
\textbf{.csv}, \textbf{.xls}, or \textbf{.xlsx}. On the other hand, DOM
parsing is carried out through programming with scripts and requires
considerable effort. Additionally, some programs or plugins facilitate
the retrieval of information without an extensive knowledge of programming.

\subsubsection{Data from social media, the API
concept}\label{data-from-social-media-the-api-concept}

In the land of social media, there have been winners and losers. Some social-media platforms, once ubiquitous, have now turned into little more than
websites, such as Myspace, or were entirely discontinued by their
creators, like Orkut and Google+. A few, however, have stood the test of
time: X (formerly Twitter), one of the most scraped social-media
platforms, Facebook, LinkedIn, Reddit, are just a few of these. Others are new or
emerging, such as TikTok. No matter their current status, their content
can be a valuable source of data. Let us briefly
consider them and the collection of their data through an API.

Twitter's name was inspired by how birds communicate, with short,
high-pitched calls or sounds directed to an audience which understands
those messages. The titular company debuted in 2006. Now rebranded as X,
it is well known for its micro-blogging consisting of short and concise messages (Murthy 2018). This model of
communication is both efficient and information-dense. Messages, called
tweets, have several key elements: \textbf{addresses}
(usernames marked by @), \textbf{text} (limited to \textbf{280}
characters), \textbf{hashtags} (marked by \textbf{\#} to denote topics),
and types of interaction like \textbf{replies}, \textbf{mentions}, and
\textbf{retweets}.

Data from X and other social media are particularly interesting for
marketing analysts because it has been, and continues to be,
used extensively for viral advertising (Tellis et al. 2019). It has strong user
engagement, a solid user base, and supports many languages. Moreover, it
contains vast amounts of real-time data. The mix of user behavior and
competitor strategy is crucial for understanding and analyzing marketing. It is an
ideal place to monitor influencers' ideas and opinions, understand why they are followed, and use this information for marketing purposes: for example, knowing how
your brand is perceived on this social-media platform. Learning this can
help you stay aware of potential discontent. Understanding how people
``feel'' about your brand and how this sentiment changes in response to
your actions is vital. This type of analysis, known as ``sentiment
analysis'', will be discussed in more detail in
Chapter~\ref{sec-sentiment-analysis}. Some R packages that
help you scrape data from X through APIs are \texttt{rtweet}
(Kearney 2019), \texttt{streamR} (Barbera 2018), and
\texttt{academictwitteR} (Barrie and Ho 2021).

With its massive user base, Facebook has been seen as an ideal platform for
promoting products, ideas, or projects. The connection between
users and the ''like'' system made it an easy way to collect personal data for free. These can be harnessed for
marketing analysis, especially for segmentation, where targeting by
geographical regions can enhance communication and promote specific
products. Two important R packages for dealing with data on Facebook
are: \texttt{Rfacebook} (Barbera et al. 2017) and \texttt{fbRads}
(Daroczi, Gopal, and Nagy 2023).

LinkedIn, a popular social-media platform for professionals, was set up to aid networking. Finding a job through LinkedIn has become a viable
option, and the personal pages can contain a lot of interesting
information, mainly about what occupies most of our time: work.
This can be extracted to check for who might be important partners or
influencers in marketing campaigns. Other analyses could be done to
understand market dynamics, emerging trends, and information about
competitors, including their workforce size or skill sets. To access
data through LinkedIn's APIs via R, you can use \texttt{Rlinkedin}
(Piccirilli 2016).

Reddit, a very popular platform for sharing opinions and asking
questions about products, recipes, medicines, supplements, and much
more, has become a rich source of data. The impressive quantity of data and
opinions make this social-media platform extremely useful for
marketing analysis. The R package to interact with this platform is \texttt{RedditExtractoR} (Rivera 2023).

All the platforms mentioned above offer, or offered, the opportunity for researchers to scrape their
data through an API. An API consists
of definitions and communication protocols that allow a person with a
computer to connect to a server providing this technology. It is a
privileged gateway to data that needs to be first identified and then
transferred to your computer. APIs generally use a
username and password system, plus or minus other related values, also known
as keys. These alphanumeric sequences are generated for a specific user and work
like credentials. APIs are interfaces; they open a door through which the user can
access data directly from the service's backend. However, it is rare that you are able to do this immediately, and collecting data through APIs can be
problematic. Significant issues arise due to constant changes in
authorization modes, data access limitations, and managerial decisions
which can restrict access. Several R packages written for this purpose
become obsolete after a few years, or even months, due to internal changes
in social-media data-accessibility policies and parameters. As such,
although scraping with APIs might seem a good way forward, other methods of
obtaining data may be more suitable for the task you have in front of you. Remember, while
working with R packages that do API scraping, two components must be
stable: the politics of the API provider and the R package.

\subsection{Scraping with your
browser}\label{scraping-with-your-browser}

Having tools that simplify the data-collection process without requiring
deep technical know-how can be useful and comforting. Take the Web
Scraper plugin for Google Chrome, for example. It is a user-friendly
tool for web scraping, requiring only a basic understanding of HTML and
CSS. You can install it easily from the Chrome store, and, as with most
plugins, a browser restart is needed post-installation. Web Scraper stands out for its:

\begin{itemize}
\tightlist
\item
  Ability to navigate and scrape multiple pages with a defined
  structure, like
  ``\href{http://www.example.com/pages/\%5B1-10\%5D\%22.}{http://www.example.com/pages/{[}1-10{]}``.}\\
\item
  Versatility in scraping different data types.\\
\item
  Ability to extract dynamically generated text from client-side
  scripts, including AJAX and JavaScript.\\
\item
  The possibility of reviewing and exporting the local storage of scraped data to a
  CSV file.\\
\item
  The ease in switching between exportable and importable sitemaps
  which detail the scope of a scraping project.
\end{itemize}

Setting it up involves defining the sitemap (which pages to visit) and the
data to be scraped. During the scraping process, Chrome will display a
series of pop-up windows showing its progress. You will receive a
notification once the scraping is complete. All the data can be
browsed and exported as a \textbf{.csv} file, ready to be imported into
R for further analysis. The plugin also allows you to export the sitemap for
future use. Check out the plugin's official webpage for a comprehensive
guide on using this tool. To see the scraping results in action, visit this
book's accompanying website.

\subsubsection{Other Plugins}\label{other-plugins}

The world of Chrome plugins for data scraping is vast and varied. Take
Instant Data Scraper, for example, a plugin that simplifies the scraping
process, requiring minimal technical know-how. It uses AI analysis to see
what data can be extracted from an HTML page, bringing something unique
to your scraper's toolbox. And for those interested in specific sites,
Amazon Reviews Exporter and Amazon Data Scraper are  plug-ins designed
for collecting data from Amazon. Wherever the data are, it is lkely there is
a tool ready to help you harvest them. This is true for a range of sites
that are well-known to be rich in data, from Parlor to Gab, and Yelp to
Trustpilot. And it is not just Chrome that offers these tools, Data Miner, for instance, works with both
Chrome and Microsoft Edge.

Focusing on something more practical, using Amazon Reviews Exporter's
free version, you can download about 10 reviews at a time (be
aware that these parameters might change over time). For example, we could grab
reviews for ``The Lord of the Rings Hardcover - Box set, March 3, 1988''
from Amazon, save them as a CSV file, and then dive into the data
using R.

\emph{Program code starts}

\begin{Shaded}
\begin{Highlighting}[]
\CommentTok{\# Set global options to limit text width}

\CommentTok{\# Load tibble package}
\FunctionTok{library}\NormalTok{(tibble)}

\CommentTok{\# Read CSV as dataframe}
\NormalTok{tolkien }\OtherTok{\textless{}{-}} \FunctionTok{read.csv}\NormalTok{(}\StringTok{"data/LordOR\_amzn.csv"}\NormalTok{, }\AttributeTok{header =} \ConstantTok{TRUE}\NormalTok{, }\AttributeTok{stringsAsFactors =} \ConstantTok{FALSE}\NormalTok{)}

\CommentTok{\# Convert dataframe to tibble for readability}
\NormalTok{tolkien\_t }\OtherTok{\textless{}{-}} \FunctionTok{as\_tibble}\NormalTok{(tolkien)}

\CommentTok{\# Display tibble structure}
\FunctionTok{str}\NormalTok{(tolkien\_t)}
\end{Highlighting}
\end{Shaded}

\begin{verbatim}
tibble [50 x 16] (S3: tbl_df/tbl/data.frame)
 $ id             : chr [1:50] "R177MIRDIAB6HL" "R1AD5LS9CM4L0R" "R1PI73EA3XE8F0" "R1LU51BWLGZ07S" ...
 $ username       : chr [1:50] "Karen Reynolds" "Julian F Waldby" "ScifiEtronica" "Matt" ...
 $ profileURL     : chr [1:50] "https://www.amazon.com/gp/profile/amzn1.account.AE7EARUOMHKHUILDESDXVYENQNUQ/ref=cm_cr_arp_d_gw_btm?ie=UTF8" "https://www.amazon.com/gp/profile/amzn1.account.AGKZU2BT3FHOUHWKGATWIWVMBS4Q/ref=cm_cr_arp_d_gw_btm?ie=UTF8" "https://www.amazon.com/gp/profile/amzn1.account.AHH32XFUTXOW33GVBF3FW6RYMZ5A/ref=cm_cr_arp_d_gw_btm?ie=UTF8" "https://www.amazon.com/gp/profile/amzn1.account.AHHRZF6OXT5UYFW44TPWFKXUUOPQ/ref=cm_cr_arp_d_gw_btm?ie=UTF8" ...
 $ rating         : int [1:50] 5 5 5 5 5 5 5 5 5 5 ...
 $ title          : logi [1:50] NA NA NA NA NA NA ...
 $ review_url     : chr [1:50] "https://www.amazon.com/gp/customer-reviews/R177MIRDIAB6HL/ref=cm_cr_arp_d_rvw_ttl?ie=UTF8&ASIN=0395489326" "https://www.amazon.com/gp/customer-reviews/R1AD5LS9CM4L0R/ref=cm_cr_arp_d_rvw_ttl?ie=UTF8&ASIN=0395489326" "https://www.amazon.com/gp/customer-reviews/R1PI73EA3XE8F0/ref=cm_cr_arp_d_rvw_ttl?ie=UTF8&ASIN=0395489326" "https://www.amazon.com/gp/customer-reviews/R1LU51BWLGZ07S/ref=cm_cr_arp_d_rvw_ttl?ie=UTF8&ASIN=0395489326" ...
 $ asin           : int [1:50] 395489326 395489326 395489326 395489326 395489326 395489326 395489326 395489326 395489326 395489326 ...
 $ content        : chr [1:50] "Such a nice looking collection. So nice that it can be used as room decor when not being read. Hey its Lord of "| __truncated__ "All the books together.  Gave me something to do whenever I got depressed over holidays." "I was looking for an edition in three volumes that included all the corrections of mistakes made in prior publi"| __truncated__ "This is a review of the hardcover boxed set of The Lord of the Rings, which includes the three volumes of the o"| __truncated__ ...
 $ date           : chr [1:50] "Reviewed in the United States on January 2, 2024" "Reviewed in the United States on January 4, 2024" "Reviewed in the United States on November 11, 2023" "Reviewed in the United States on April 21, 2021" ...
 $ format_strip   : logi [1:50] NA NA NA NA NA NA ...
 $ verified       : chr [1:50] "Verified Purchase" "Verified Purchase" "Verified Purchase" "Verified Purchase" ...
 $ helpful        : chr [1:50] "One person found this helpful" "" "14 people found this helpful" "112 people found this helpful" ...
 $ videoURL       : logi [1:50] NA NA NA NA NA NA ...
 $ video_cover_url: logi [1:50] NA NA NA NA NA NA ...
 $ imgs           : chr [1:50] "[]" "[]" "[]" "[\"https://images-na.ssl-images-amazon.com/images/W/MEDIAX_792452-T2/images/I/811XmQymqDL._SY88.jpg\",\"https:/"| __truncated__ ...
 $ normalImgs     : chr [1:50] "[]" "[]" "[]" "[\"https://images-na.ssl-images-amazon.com/images/W/MEDIAX_792452-T2/images/I/811XmQymqDL.jpg\",\"https://image"| __truncated__ ...
\end{verbatim}

\begin{Shaded}
\begin{Highlighting}[]
\CommentTok{\# Show first 2 entries}
\FunctionTok{head}\NormalTok{(tolkien\_t}\SpecialCharTok{$}\NormalTok{content, }\DecValTok{2}\NormalTok{)}
\end{Highlighting}
\end{Shaded}

\begin{verbatim}
[1] "Such a nice looking collection. So nice that it can be used as room"    
[2] "decor when not being read. Hey its Lord of the Rings and a perfect set."
[3] "All the books together.  Gave me something to do whenever I got"        
[4] "depressed over holidays."                                               
\end{verbatim}

\emph{Program code ends}

As you can see, even from a small dataset, as in the one above, you can
extract a lot of information and save it in a structured form. And,
the more data you have access to, the better and more accurate the analysis will be. All in
all, there there is no one size fits all in the scraping-plugins world.
Some will work better than others, some new tools will appear on the market and others will cease to exist. But it is important that you
know, as a data analyst, that you can always find another plugin for your data-scraper tool box. Most of these services are not free or have a free version that
allows you to collect a limited number of pages or items. Alternatively,
if you need to have full control over the downloaded content, you can be
creative and use your programming skills to collect information, for which, see
next section.

\subsection{Scraping by programming}\label{scraping-by-programming}

To the non-specialist, a webpage appears as a blend of text, images,
and videos, all presented on a web browser. However, beneath
this surface lies a more complex reality of a double-language presence
made of a markup language and a style sheet that instructs the browser on what
and how to display the content. The structure of the page is governed
by hyper text markup language (HTML), while its visual presentation is
run by cascading style sheets (CSS). HTML has an easy-to-understand
syntax. It gives directives to the web browser and embeds text in HTML
pages. The text displayed on a webpage is (nearly) always enclosed in
HTML tags such as \texttt{tag\textgreater{}}Some
Text\texttt{\textless{}tag\textgreater{}}. The text inside the angular
brackets \texttt{\textless{}\textgreater{}} can be any HTML keyword such
as \texttt{\textless{}html\textgreater{}},
\texttt{\textless{}body\textgreater{}},
\texttt{\textless{}h1\textgreater{}},
\texttt{\textless{}p\textgreater{}}, among others. The slash inside the brackets
\texttt{\textless{}/\textgreater{}} signifies the end of an element. In
technical jargon, the term ``element'' defines the opening and closing
of two tags and the content between them. Once a tag is open, it must
also be closed, like a bracket in mathematics. There are some
exceptions, though, where the tags are self-contained and do not require
closing, such as \texttt{\textless{}br\textgreater{}} for a line break.
Some HTML tags contain attributes, such as \texttt{id} or
\texttt{class}. They link HTML elements to the CSS for
the aesthetic styling of the webpage. An important tag in HTML is
\texttt{\textless{}div\textgreater{}}, used in web design and layout, and is essential for structuring the site. It is commonly used together with CSS
for styling and JavaScript for functionality, so pay attention to it
whenever you find it. Figure~\ref{fig-paragraphelement} shows the main
components of an HTML element.

\begin{figure}

\centering{

\includegraphics[width=0.9\textwidth,height=\textheight]{figure/chapter3/paragraphelement.png}

}

\caption{\label{fig-paragraphelement}HTML, paragraph element}

\end{figure}%

A simplified example of an HTML code, that also includes an element like that
in Figure~\ref{fig-paragraphelement}, can be seen in the snippet
below:

\emph{Programming code starts}

\begin{Shaded}
\begin{Highlighting}[]
   \DataTypeTok{\textless{}!DOCTYPE}\NormalTok{ html}\DataTypeTok{\textgreater{}}
\DataTypeTok{\textless{}}\KeywordTok{html}\DataTypeTok{\textgreater{}}
\DataTypeTok{\textless{}}\KeywordTok{head}\DataTypeTok{\textgreater{}}
    \DataTypeTok{\textless{}}\KeywordTok{title}\DataTypeTok{\textgreater{}}\NormalTok{Book}\DataTypeTok{\textless{}/}\KeywordTok{title}\DataTypeTok{\textgreater{}}
    \DataTypeTok{\textless{}}\KeywordTok{meta}\OtherTok{ name}\OperatorTok{=}\StringTok{"description"}\OtherTok{ content}\OperatorTok{=}\StringTok{"Marketing"}\DataTypeTok{\textgreater{}}
    \DataTypeTok{\textless{}}\KeywordTok{script}\OtherTok{ type=}\StringTok{"text/javascript"}\DataTypeTok{\textgreater{}} 
    \CommentTok{// JavaScript code can go here}
    \DataTypeTok{\textless{}/}\KeywordTok{script}\DataTypeTok{\textgreater{}}
    \DataTypeTok{\textless{}}\KeywordTok{link}\OtherTok{ rel}\OperatorTok{=}\StringTok{"stylesheet"}\OtherTok{ href}\OperatorTok{=}\StringTok{"style.css"}\DataTypeTok{\textgreater{}}
\DataTypeTok{\textless{}/}\KeywordTok{head}\DataTypeTok{\textgreater{}}
\DataTypeTok{\textless{}}\KeywordTok{body}\DataTypeTok{\textgreater{}}
    \DataTypeTok{\textless{}}\KeywordTok{h1}\DataTypeTok{\textgreater{}}\NormalTok{Scraping}\DataTypeTok{\textless{}/}\KeywordTok{h1}\DataTypeTok{\textgreater{}}
    \DataTypeTok{\textless{}}\KeywordTok{article}\DataTypeTok{\textgreater{}}
        \DataTypeTok{\textless{}}\KeywordTok{p}\DataTypeTok{\textgreater{}}
\NormalTok{            Paragraph with }\DataTypeTok{\textless{}}\KeywordTok{sub}\DataTypeTok{\textgreater{}}\NormalTok{subscript}\DataTypeTok{\textless{}/}\KeywordTok{sub}\DataTypeTok{\textgreater{}}\NormalTok{ and }\DataTypeTok{\textless{}}\KeywordTok{sup}\DataTypeTok{\textgreater{}}\NormalTok{superscript}\DataTypeTok{\textless{}/}\KeywordTok{sup}\DataTypeTok{\textgreater{}}\NormalTok{ text.}
        \DataTypeTok{\textless{}/}\KeywordTok{p}\DataTypeTok{\textgreater{}}
        \DataTypeTok{\textless{}}\KeywordTok{p}\OtherTok{ id}\OperatorTok{=}\StringTok{"par2"}\DataTypeTok{\textgreater{}}
\NormalTok{            Paragraph with }\DataTypeTok{\textless{}}\KeywordTok{i}\DataTypeTok{\textgreater{}}\NormalTok{italic}\DataTypeTok{\textless{}/}\KeywordTok{i}\DataTypeTok{\textgreater{}}\NormalTok{ and }\DataTypeTok{\textless{}}\KeywordTok{b}\DataTypeTok{\textgreater{}}\NormalTok{bold}\DataTypeTok{\textless{}/}\KeywordTok{b}\DataTypeTok{\textgreater{}}\NormalTok{ text.}
        \DataTypeTok{\textless{}/}\KeywordTok{p}\DataTypeTok{\textgreater{}}
    \DataTypeTok{\textless{}/}\KeywordTok{article}\DataTypeTok{\textgreater{}}
    \DataTypeTok{\textless{}}\KeywordTok{table}\DataTypeTok{\textgreater{}}
\DataTypeTok{\textless{}}\KeywordTok{tr}\DataTypeTok{\textgreater{}}
\DataTypeTok{\textless{}}\KeywordTok{td}\DataTypeTok{\textgreater{}}\NormalTok{Tablets}\DataTypeTok{\textless{}/}\KeywordTok{td}\DataTypeTok{\textgreater{}}
\DataTypeTok{\textless{}/}\KeywordTok{tr}\DataTypeTok{\textgreater{}}
\DataTypeTok{\textless{}/}\KeywordTok{table}\DataTypeTok{\textgreater{}}
    \DataTypeTok{\textless{}}\KeywordTok{ul}\DataTypeTok{\textgreater{}}
        \DataTypeTok{\textless{}}\KeywordTok{li}\DataTypeTok{\textgreater{}}\NormalTok{API}\DataTypeTok{\textless{}/}\KeywordTok{li}\DataTypeTok{\textgreater{}}
        \DataTypeTok{\textless{}}\KeywordTok{li}\DataTypeTok{\textgreater{}}\NormalTok{DOM}\DataTypeTok{\textless{}/}\KeywordTok{li}\DataTypeTok{\textgreater{}}
    \DataTypeTok{\textless{}/}\KeywordTok{ul}\DataTypeTok{\textgreater{}}
    \DataTypeTok{\textless{}}\KeywordTok{footer}\DataTypeTok{\textgreater{}}
        \DataTypeTok{\textless{}}\KeywordTok{p}\DataTypeTok{\textgreater{}}\NormalTok{This is a footer paragraph.}\DataTypeTok{\textless{}/}\KeywordTok{p}\DataTypeTok{\textgreater{}}
        \DataTypeTok{\textless{}}\KeywordTok{a}\OtherTok{ href}\OperatorTok{=}\StringTok{"https://www.danieldanlab.com"}\OtherTok{Link to a page}\ErrorTok{\textless{}/}\OtherTok{a}\DataTypeTok{\textgreater{}}
    \DataTypeTok{\textless{}/}\KeywordTok{footer}\DataTypeTok{\textgreater{}}
\DataTypeTok{\textless{}/}\KeywordTok{body}\DataTypeTok{\textgreater{}}
\DataTypeTok{\textless{}/}\KeywordTok{html}\DataTypeTok{\textgreater{}}
\end{Highlighting}
\end{Shaded}

\emph{Programming code ends}

Using programming languages, for example, R, Python, or Java, we can
navigate through the internal elements of these web pages and retrieve valuable
information. This process involves parsing the pages into a document-object-model (DOM) tree structure (Keith and Sambells 2011). Let us
demystify the link between HTML and DOM. When a
browser loads a webpage, it parses the web page sequentially, reading
the HTML tags in sequence. While parsing, it builds the DOM, and the
HTML tags are transformed into DOM nodes. If, during the parsing
process, there are scripts, JavaScripts, for example, or external
elements such as CSS directives and images that might interrupt the
process, the DOM building still goes ahead. Most of the time, in
parallel, with little or no interruptions. So, the DOM represents a
model of a page, while HTML is a language that allows you to create the
DOM. After its creation, scripts can manipulate it to change a page's
content, structure, or style after it has been loaded. This procedure
allows for it to be interactive, making pages dynamic and accessible alike. It is
easier to locate specific data on a webpage through the hierarchical
tree structure of the DOM, and scraping tools can use the DOM to identify
\texttt{ids}, \texttt{classes}, and \texttt{tag\ names} to extract the information you need. For a visual representation of how a basic HTML
page, like that in the previous example, is organized by the DOM, refer to
Figure~\ref{fig-DOMTree}. Be aware that only the opening
tags are represented, the closing section is, for the most part, a mirror of these.

\begin{figure}

\centering{

\includegraphics[width=0.9\textwidth,height=\textheight]{figure/chapter3/html_2.png}

}

\caption{\label{fig-DOMTree}DOM Structure}

\end{figure}%

The \texttt{\textless{}html\textgreater{}} node is the \textbf{root} of
the tree, it has two branches: the \texttt{\textless{}head\textgreater{}}
and \texttt{\textless{}body\textgreater{}}. It tells the browser that it
encapsulates all the  elements of the page. The
\texttt{\textless{}head\textgreater{}} node is like the preface of a book. As well as other elements, it
provides meta-information about what lies underneath the page.
The \texttt{\textless{}meta\textgreater{}} tag adds information about
the HTML page. It does not have a closing counterpart.
\texttt{\textless{}title\textgreater{}},
\texttt{\textless{}script\textgreater{}}, and
\texttt{\textless{}link\textgreater{}} are just some of the tags you might encounter in a header. The \texttt{\textless{}body\textgreater{}} node
contains the main content, structured into headers, such as
\texttt{\textless{}h1\textgreater{}}, the main content, which can be made up of
articles \texttt{\textless{}article\textgreater{}}, tables,
\texttt{\textless{}table\textgreater{}}, and unordered lists
\texttt{\textless{}ul\textgreater{}}, and ends with the footer
\texttt{\textless{}footer\textgreater{}}.
\texttt{\textless{}h1\textgreater{}} is the heading tag with the
largest dimensions. If you want to create smaller headings,
\texttt{\textless{}h2\textgreater{}} to
\texttt{\textless{}h6\textgreater{}} offer progressively smaller sizes.
\texttt{\textless{}li\textgreater{}...\textless{}/li\textgreater{}}
show that we will enclose a \textbf{li}st between the tags. To make an
\textbf{u}nordered \textbf{l}ist, you use the
\texttt{\textless{}ul\textgreater{}} tag. Of course, there is a tag for
an \textbf{o}rdered \textbf{l}ist too which adds numbers or letters in
front of \texttt{\textless{}ol\textgreater{}}. You can create
a table with a unique element that has a table row
\texttt{\textless{}tr\textgreater{}}, within which is a \textbf{t}able
\textbf{c}ell not defined by the
\texttt{\textless{}tc\textgreater{}} tag, but
\texttt{\textless{}td\textgreater{}}, which stands for \textbf{t}able
\textbf{d}ata (or cell). For more information on HTML and CSS, check the
additional material on the https://danieldanlab.github.io webpage.

In the following example, we have imported a webpage into R and removed the
unnecessary part by stripping away the code and liberating the embedded
text.

The \texttt{rvest} package (Wickham 2022a) provides functions that allow
us to scrape (or ha\textbf{rvest}) data by parsing the HTML content of
the webpage. \texttt{rvest} integrates with the
\texttt{magrittr} package seamlessly, using the pipe operator
\texttt{\%\textgreater{}\%} to pass output easily from one function to
the next. This approach is detailed in the introductory chapter, but let us
explore how things work in practice with an example.

\emph{Program code starts}

\begin{Shaded}
\begin{Highlighting}[]
\CommentTok{\# Load necessary library}
\FunctionTok{library}\NormalTok{(rvest)}

\CommentTok{\# Define the URL to scrape}
\NormalTok{url }\OtherTok{\textless{}{-}} \StringTok{"https://webscraper.io/test{-}sites/e{-}commerce/allinone"}

\CommentTok{\# Read and parse the HTML from the URL}
\NormalTok{parsed\_html }\OtherTok{\textless{}{-}} \FunctionTok{read\_html}\NormalTok{(url)}

\CommentTok{\# Extract product titles}
\NormalTok{titles }\OtherTok{\textless{}{-}}\NormalTok{ parsed\_html }\SpecialCharTok{\%\textgreater{}\%}
  \FunctionTok{html\_nodes}\NormalTok{(}\StringTok{"h4 .title"}\NormalTok{) }\SpecialCharTok{\%\textgreater{}\%}
  \FunctionTok{html\_text}\NormalTok{()}

\CommentTok{\# Extract product descriptions}
\NormalTok{descriptions }\OtherTok{\textless{}{-}}\NormalTok{ parsed\_html }\SpecialCharTok{\%\textgreater{}\%}
  \FunctionTok{html\_nodes}\NormalTok{(}\StringTok{".description"}\NormalTok{) }\SpecialCharTok{\%\textgreater{}\%}
  \FunctionTok{html\_text}\NormalTok{()}

\CommentTok{\# Extract product prices as text}
\NormalTok{prices\_text }\OtherTok{\textless{}{-}}\NormalTok{ parsed\_html }\SpecialCharTok{\%\textgreater{}\%}
  \FunctionTok{html\_nodes}\NormalTok{(}\StringTok{".caption .pull{-}right"}\NormalTok{) }\SpecialCharTok{\%\textgreater{}\%}
  \FunctionTok{html\_text}\NormalTok{()}

\CommentTok{\# Clean and convert prices to numeric}
\NormalTok{prices }\OtherTok{\textless{}{-}} \FunctionTok{as.numeric}\NormalTok{(}\FunctionTok{gsub}\NormalTok{(}\StringTok{"}\SpecialCharTok{\textbackslash{}\textbackslash{}}\StringTok{$"}\NormalTok{, }\StringTok{""}\NormalTok{, prices\_text))}

\CommentTok{\# Combine extracted data into a data frame}
\NormalTok{product\_data }\OtherTok{\textless{}{-}} \FunctionTok{data.frame}\NormalTok{(titles, descriptions, prices)}

\CommentTok{\# Preview the first few rows of the data frame}
\FunctionTok{head}\NormalTok{(product\_data)}
\end{Highlighting}
\end{Shaded}

\emph{Program code ends}

The \texttt{read\_html()} function is designed to parse the content of
an HTML file. Once the file is parsed and stored in a variable, such as
\texttt{titles}, we need to pass this variable containing the entire
structure of the page and the text to the \texttt{html\_nodes()}
function. This will extract the text using various HTML
and CSS selectors, in this case, \texttt{h4} and \texttt{.title}. If you
were wondering why we use more than one selector, it is because this allows us to be specific and pinpoint the content we
need to extract. It is important to be able to identify these selectors visually. An easy
way to access and inspect them is through the developer's web tools. All
modern browsers can show a page's source and unveil its DOM. On some
browsers, in the settings, which are usually represented by three vertical dots or
lines, you need to select \textbf{More tools}, then \textbf{Developer
tools}. A menu for the developer tools will appear either at the bottom
of the page or on the right-hand side. The first step is to select an element
on the page. Let us identify the \textbf{title} of the first item whose
data we want to have downloaded. Given the repetitive nature of the code
structure, identifying the pattern for one item is usually sufficient
as the other items will likely follow the same pattern. Click on the
icon depicted in Figure~\ref{fig-icon} and then click on the title of
one product, see Figure~\ref{fig-product}. Usually, such a procedure is
applied to all the elements to be harvested from a page, regardless of
the number of products. The code above is sufficient to harvest all the features you want of all the products on a page. Be aware that what you
see in the developer tools of a browser is the DOM, not just the raw
HTML, as it displays the current state of the DOM as it is interpreted
and rendered by the browser being used.

\begin{figure}

\centering{


\includegraphics[width=0.1\textwidth,height=\textheight]{figure/chapter3/selector.png}

}

\caption{\label{fig-icon}Icon to select elements on a webpage, image
from a Chrome browser, compiled by the authors}

\end{figure}%

\begin{figure}

\centering{

\includegraphics[width=0.3\textwidth,height=\textheight]{figure/chapter3/product.png}

}

\caption{\label{fig-product}Title selection of a product, image from a
web page, compiled by the authors}

\end{figure}%

The HTML snippet code
will be displayed in the \textbf{Developer tools} frame underneath,  allowing you to detect the path where the title is
situated, see Figure~\ref{fig-snippet}.

\begin{figure}

\centering{

\includegraphics[width=0.9\textwidth,height=\textheight]{figure/chapter3/snippet.png}

}

\caption{\label{fig-snippet}Code snippet underlining the path to follow
to obtain the title's text}

\end{figure}%

When examining the code, you can build an imaginary path going through
both elements, \texttt{h4} and \texttt{.title}, to locate a specific
title, such as ``Acer Aspire 3 A315-15 Black''. The function is designed
to extract a node containing text but does not know the path to follow.
This is why we have to inspect the code and supply the function with the right elements. Essentially, it selects only the desired node,
removing all the unnecessary information. This output is then passed to
the \texttt{html\_text()} function, which extracts the text from
the node. This approach is applied in the same way to all the other variables.
Note that the \textbf{prices\_text} variable has a character sign before
the numerical value. With the help of a regular expression you have
studied, {[}here{]}, we can remove that sign and convert the remaining
string made by character digits into an integer. This step is always
needed for subsequent mathematical computations: mathematical operations
are done using numbers, not characters.

\subsection{Scraping by programming +}\label{scraping-by-programming-1}

Sometimes, scraping web pages based on a static pattern is very
challenging, especially when actions like clicking a link or navigating
to a ``next'' page are required. In such cases, we need a tool which can
mimic and automate human actions on a webpage. Selenium (2024) has a set
of various open-source software tools, each with a distinct
method, to help do this. Primarily used for automatizing tests concerning the
functionality of websites, Selenium can also be used as a remote control. It
makes it possible for the interface elements on web pages, for example,
links or pop-ups, to be accessed via scripting rather than through clicking with a mouse.
It can even handle JavaScript processes as well (Garcı́a et al. 2020).
The procedures mentioned above are usually tricky to execute without
Selenium due to the problem of integrating JavaScript with the browser.

Another component is Selenium's ability to identify the agent,
essentially the browser used for scraping. Common choices include Google
Chrome (the default option), Firefox, or Edge, and these can be used in
their standard or headless versions. A headless browser provides
automated web-page interaction without a graphical user interface. It
can be chosen in lieu of commonly used web browsers to minimize
interactions, and so maximizing the user's anonymity. The key packages used for
this task are \texttt{RSelenium} (Harrison 2022) and \texttt{rvest}. The
former ensures that R and Selenium interact and provides a set
of R bindings from the Selenium WebDriver (Garcı́a et al. 2021), while
the latter handles the scraping. The script starts the
Selenium server sequentially and pauses the system for two seconds, allowing it to be
loaded. Then, it links the driver to the browser, which opens a session and loads the external file. Eventually, the data is
stored in a data frame and then saved in CSV format. However, be aware
of potential compatibility issues between Chrome and the versions of its driver.
In such cases, switching to Firefox or another browser might resolve the
issue. Also, if a script fails due to a communication port being busy,
consider closing the current session or choosing a different
port. For more detailed guidance, visit our website.

Let us use this tool on a page we are already familiar with from the
webscraper.io website, starting from a general page. The function
\texttt{rsDriver()} may default to using the Chrome browser if not
specified by the user or the user's environment. The option
\texttt{verbose\ =\ FALSE} suppresses the system messages. Its default
value is TRUE. We suggest experimenting with the code and its various
attributes in order to understand better the true potential of these
functions.

\emph{Program code starts}

\begin{Shaded}
\begin{Highlighting}[]
\CommentTok{\# Load RSelenium}
\FunctionTok{library}\NormalTok{(RSelenium)}

\CommentTok{\# Set URL}
\NormalTok{url }\OtherTok{\textless{}{-}} \StringTok{"https://webscraper.io/test{-}sites/e{-}commerce/allinone"}

\CommentTok{\# Initialize Selenium driver}
\NormalTok{driver }\OtherTok{\textless{}{-}} \FunctionTok{rsDriver}\NormalTok{(}\AttributeTok{verbose =} \ConstantTok{FALSE}\NormalTok{)}

\CommentTok{\# Create client}
\NormalTok{remote\_driver }\OtherTok{\textless{}{-}}\NormalTok{ driver}\SpecialCharTok{$}\NormalTok{client}

\CommentTok{\# Navigate to URL}
\NormalTok{remote\_driver}\SpecialCharTok{$}\FunctionTok{navigate}\NormalTok{(url)}

\CommentTok{\# Wait for page load}
\FunctionTok{Sys.sleep}\NormalTok{(}\DecValTok{2}\NormalTok{)}

\CommentTok{\# Click element by XPath}
\NormalTok{button\_element }\OtherTok{\textless{}{-}}\NormalTok{ remote\_driver}\SpecialCharTok{$}\FunctionTok{findElement}\NormalTok{(}\AttributeTok{using =} \StringTok{"xpath"}\NormalTok{, }
                    \AttributeTok{value =} \StringTok{\textquotesingle{}//*[@id="side{-}menu"]/li[3]/a\textquotesingle{}}\NormalTok{)}
\NormalTok{button\_element}\SpecialCharTok{$}\FunctionTok{clickElement}\NormalTok{()}

\CommentTok{\# Get and print new URL}
\NormalTok{url1 }\OtherTok{\textless{}{-}}\NormalTok{ remote\_driver}\SpecialCharTok{$}\FunctionTok{getCurrentUrl}\NormalTok{()}
\FunctionTok{print}\NormalTok{(url1)}

\CommentTok{\# Scrape data (code for scraping goes here)}

\CommentTok{\# Close driver and stop server}
\NormalTok{remote\_driver}\SpecialCharTok{$}\FunctionTok{close}\NormalTok{()}
\NormalTok{driver}\SpecialCharTok{$}\NormalTok{server}\SpecialCharTok{$}\FunctionTok{stop}\NormalTok{()}
\end{Highlighting}
\end{Shaded}

\emph{Program code ends}

To sum up, the code starts by loading the necessary library and
setting a target URL. It then identifies and follows a path to a
specific element on the webpage, simulating a click. This action leads
us to a new page with an updated URL. Effectively, this code automates the process of navigating through web pages, enabling important
information to be scraped without the need for manual interaction, such as clicking on
links.

\subsection{Data from databases: SQL,
NoSQL}\label{data-from-databases-sql-nosql}

Structured Query Language (SQL) is a programming language designed for
accessing and managing databases, regardless of their physical location.
It is used to retrieve, write, or update inherently structured data which
are stored in tables linked through a relational scheme. As
has been highlighted above, structured data is the preferred format for
analysts and SQL has a distinctive syntax, which can be used to interact
directly with databases. While direct SQL commands can be complex,
wrapper functions are available to ``conceal'' the SQL jargon and
simplify these commands into more user-friendly terms. To illustrate,
consider a hypothetical SQL query: In the \texttt{users} table, for the
\texttt{user\_id} identified by number 25, select columns like
\texttt{name}, \texttt{familyname}, and \texttt{address}. This example,
though not connected to an actual database, demonstrates the fundamental
power of SQL for manipulating data.

\begin{Shaded}
\begin{Highlighting}[]
\KeywordTok{SELECT}\NormalTok{ name, familyname, address }\KeywordTok{FROM}\NormalTok{ users }\KeywordTok{WHERE}\NormalTok{ user\_id }\OperatorTok{=} \DecValTok{25} 
\end{Highlighting}
\end{Shaded}

If we replace \texttt{name}, \texttt{familyname}, and \texttt{address}
with an asterisk (*) in the SQL query, it will select all the fields from
the table. Additionally, omitting the WHERE clause will result in the
selection of all users, not just those matching a specific condition.

\begin{Shaded}
\begin{Highlighting}[]
\KeywordTok{SELECT} \OperatorTok{*} \KeywordTok{FROM}\NormalTok{ users}
\end{Highlighting}
\end{Shaded}

The indisputable strength of R lies in its built-in functions, which
enable the quick analysis of data and the ability to visualize results through
graphs for immediate comprehension. Moreover, R, with its extensive
range of packages, offers drivers for connecting to various types of
databases, including MySQL, ORACLE, PostgreSQL, Impala, Amazon RedShift,
and Spark, to name just a few.

You can either use the SQL syntax embedded in simple R code or you can
use packages like \texttt{dplyr}. For the first
approach, you might use a command like:

\begin{Shaded}
\begin{Highlighting}[]
\FunctionTok{tbl}\NormalTok{(mydatabase, }\FunctionTok{sql}\NormalTok{(}\StringTok{"SELECT name, familyname, address FROM users WHERE }
\StringTok{                    user\_id = 25"}\NormalTok{))}
\end{Highlighting}
\end{Shaded}

For the second:

\begin{Shaded}
\begin{Highlighting}[]
\NormalTok{myselection }\OtherTok{\textless{}{-}} \FunctionTok{tbl}\NormalTok{(mydatabase, }\StringTok{"users"}\NormalTok{)}
\NormalTok{myselection }\SpecialCharTok{\%\textgreater{}\%}
  \FunctionTok{filter}\NormalTok{(}\AttributeTok{user\_id =} \DecValTok{25}\NormalTok{) }\SpecialCharTok{\%\textgreater{}\%}
  \FunctionTok{select}\NormalTok{(name, familyname, address)}
\end{Highlighting}
\end{Shaded}

In either case, whether using SQL directly or via \texttt{dplyr},
the richness and complexity of the dplyr package can work like a charm,
offering a vast array of analytical possibilities.

The approach to accessing data differs slightly for NoSQL databases but
retains some similarities. In NoSQL, data are often stored in JSON
format with embedded attributes. For instance, in a users table within a
\texttt{mydatabase} database, fields like \texttt{name},
\texttt{familyname}, and \texttt{address} are stored as attributes
within a single document. This helps horizontal scalability,
allowing for large amounts of data to be contained within a single
table.

This brief foray into the world of databases is intended to provide a
more comprehensive view of the diverse sources from which data can be
acquired. However, it is important to note that throughout this book, we
will primarily focus on other methods of data acquisition.

\subsection{Using services based on AI to obtain
data}\label{use-services-based-on-ai-to-obtain-data}

Because scraping using traditional software or self-produced code can be
problematic due to changing website layouts, developments in
AI---especially Large Language Models (LLMs)---offer new ways to improve
web scraping, as described in Posedaru et al. (2024). The authors propose a
solution involving a JavaScript application that uses natural language
to extract information and save the scraped results to a JSON file. The
steps include providing a URL address and description of the data that
needs to be extracted. The HTML structure is then collected, cleaned, and
transferred to ChatGPT in several chunks. ChatGPT is used to detect the
appropriate selectors corresponding to the data in the HTML structure,
after which it generates a scraper based on these elements.
Once the scraper is executed, the data is saved in a JSON file.

There are other ChatGPT-based tools, such as the gpt-automated-web-scraper,
which generates code based on user inputs and website analysis (Breeuwer
2024), or the ChatGPT Web Scraper API plugin (OpenAI 2024). These
tools enable autonomous data collection using LLMs, though ChatGPT-like
tools cannot scrape websites directly (as of now). Additionally, several
web services use AI to scrape data without the need for any coding,
such as browser.ai, bardeen.ai, and kadoa.com.

All these tools are very promising, although some will endure and others may
fade away. The important thing to understand is how to collect
data while taking into account privacy and ethics.

\section{Other sources for obtaining text data
sets}\label{other-sources-to-obtain-text-data-sets}

Gathering data for analysis, whether through web scraping or utilizing
existing datasets, can be daunting. However, readily available resources
on the web can ease this process. Notable among these are well-known
dataset aggregators like Kaggle and GitHub. These platforms host various
datasets covering various topics, including numerical data, categorical
data, and combinations thereof. In this book, we are particularly
interested in datasets suitable for natural language processing (NLP)
and analysis.

We have compiled a list of datasets from the Elite Data Science website
(\url{https://elitedatascience.com/datasets),} which are excellent for
honing your text analytics skills. We encourage you to explore these
datasets and consider their potential application for marketing:

\begin{itemize}
\tightlist
\item
  Game of Thrones: Explore the political landscape, characters, and
  battles of the popular TV series, based on George R.R. Martin's `A
  Song of Fire and Ice.'\\
\item
  World University Rankings: Analyze three global university rankings,
  each with a unique methodology and different insights.\\
\item
  IMDB 5000 Movie Dataset: Can a movie's popularity be predicted before
  its release?\\
\item
  Enron Dataset: Email data from Enron's senior management, a resource
  for understanding corporate communication. This dataset was originally
  made public and posted to the web by the Federal Energy Regulatory
  Commission during its investigation.\\
\item
  Amazon Reviews: Examine around 35 million Amazon reviews, a rich
  source for sentiment analysis and customer feedback studies.\\
\item
  Newsgroup Classification: Practice text classification and topic
  modeling with 20,000 documents from 20 newsgroups.
\end{itemize}

Check our website for more links to usable datasets.

\emph{Exercise starts}

\textbf{Exercise 1}: Inspect the Amazon dataset's first 20 and last
20 review elements. Hint: type help(head) in the R console to learn how
to do it.\\
\textbf{Exercise 2}: Make a search on the Internet and extract
unstructured data from pdf, doc, or docx files. Then, save them into a
dataframe.\\
\textbf{Exercise 3}: Choose a popular online forum or discussion board.
Extract and analyze the most recent ten posts from a subreddit
which relates to your interests. Hint: investigate the structure of the forum's URLs
and how posts are displayed to determine your scraping strategy.\\
\textbf{Exercise 4}: Find a website of your choice and use the Web
Scraper plugin to extract data.\\
\textbf{Exercise 5}: Find a website of your choice and use RSelenium and
rvest to download a usable dataset. Consider websites requiring
interaction, such as clicking buttons or filling forms, to navigate to
the data.\\
\emph{Exercise ends}

\textbf{Conclusion}

\emph{Conclusion starts}

This chapter is a resource for those beginning their journey in text
data acquisition and analysis in R. It offers practical insights into, and a
tangible foundation in, various methods of and tools for data acquisition. Text
is present in various file formats, and it has covered methods for handling PDFs, Word
documents (doc, docx), and plain text files, demonstrating
how to convert these formats into structured forms suitable for analysis
in R. APIs were another point of interest, particularly for collecting
data from social-media platforms and other online services. Web scraping
techniques were tackled, using R packages such as \texttt{rvest} and
\texttt{RSelenium} to extract data from both static and dynamic web
content, including handling pages that require user interaction. The
discussion then shifted to working with databases and how to access
structured data from both SQL and NoSQL databases. Online platforms such
as Kaggle and GitHub were mentioned as valuable sources for ready-to-use
datasets. The chapter also provides practical exercises that range from basic
data inspection to advanced web-scraping techniques. These exercises are
designed to enhance the reader's hands-on experience with a variety of
data sources and tools.

\emph{Conclusion ends}